\documentclass[11pt,letterpaper]{article}
\usepackage[utf8]{inputenc}
\usepackage[spanish]{babel}
\usepackage[margin=2.3cm]{geometry}
\usepackage{amsmath,amssymb,amsfonts}
\usepackage{graphicx}
\usepackage{longtable}
\usepackage{booktabs}
\usepackage[table]{xcolor}
\usepackage{array}
\usepackage{titlesec}
\usepackage{fancyhdr}
\usepackage{enumitem}
\usepackage{listings}
\usepackage{hyperref}
\usepackage{float}
\usepackage{tikz}
\usetikzlibrary{positioning, arrows.meta, shapes.geometric, fit, calc, matrix}

\definecolor{lightgray}{RGB}{245,245,245}
\definecolor{encabezado}{RGB}{220,220,220}
\definecolor{colorlink}{RGB}{0,0,139}

\hypersetup{colorlinks=true, linkcolor=colorlink, urlcolor=colorlink}

\lstset{
    language=SQL,
    backgroundcolor=\color{lightgray},
    basicstyle=\small\ttfamily,
    keywordstyle=\bfseries\color{colorlink},
    frame=single,
    breaklines=true,
    showstringspaces=false,
    numbers=left,
    numberstyle=\tiny
}

\titleformat{\section}{\Large\bfseries}{\thesection}{1em}{}[\titlerule]
\titleformat{\subsection}{\large\bfseries}{\thesubsection}{1em}{}
\titleformat{\subsubsection}{\normalsize\bfseries}{\thesubsubsection}{1em}{}

\pagestyle{fancy}
\fancyhf{}
\lhead{\small Plataforma DDP-IPN}
\rhead{\small Base de Datos - Documentación Técnica}
\cfoot{\thepage}

\tikzset{
    tablaUML/.style={
        matrix of nodes,
        row sep=-\pgflinewidth,
        column sep=-\pgflinewidth,
        nodes={
            rectangle, draw=black, align=left,
            font=\scriptsize\ttfamily,
            inner xsep=6pt, inner ysep=4pt,
            text width=5.8cm
        },
        column 1/.style={nodes={text width=5.8cm}},
        row 1/.style={nodes={fill=encabezado, font=\bfseries\small, align=center}}
    },
    lineaRel/.style={draw=black, thick, -{Stealth[length=6pt]}}
}

\begin{document}

\tableofcontents
\newpage


\subsection{Relación con los módulos del sistema}
\begin{longtable}{|p{4cm}|p{4.5cm}|p{6cm}|}
\hline
\rowcolor{lightgray} \textbf{Módulo} & \textbf{Tablas principales} & \textbf{Funcionalidad} \\
\hline
Quejoso & Usuarios, Quejas, Evidencias & Registro de quejas y consulta de estatus \\
\hline
Recepción & Quejas, Tutores & Validación inicial y asignación \\
\hline
Primer Contacto & Quejas, Historial & Análisis preliminar \\
\hline
Titular & Quejas, Notificaciones & Canalización de expedientes \\
\hline
Subdefensoría & Quejas, Historial & Investigación y resolución \\
\hline
\end{longtable}

\section{Modelo Entidad–Relación (DER)}

\subsection{Diagrama general}
\begin{figure}[H]
\centering
\resizebox{0.95\textwidth}{!}{
\begin{tikzpicture}[node distance=1.8cm, auto]
\matrix (gQuejas) [tablaUML] {
    |[fill=encabezado]| \textbf{Quejas} \\
    id : BIGINT (PK) \\
    usuario\_id : BIGINT (FK) \\
    descripcion : TEXT \\
    fecha\_hechos : TIMESTAMP \\
    fecha\_registro : TIMESTAMP \\
    folio : VARCHAR(255) \\
    asunto : VARCHAR(255) \\
    unidad\_donde\_ocurrio : VARCHAR(255) \\
    estatus : VARCHAR(50) \\
    nombre\_denunciado : VARCHAR(255) \\
    correo\_quejoso : VARCHAR(255) \\
    nombre\_quejoso : VARCHAR(255) \\
    identificacion\_institucional : VARCHAR(255) \\
    tipo\_identificacion : VARCHAR(255) \\
    fecha\_nacimiento : VARCHAR(255) \\
};
\matrix (gUsuarios) [tablaUML, above=2.8cm of gQuejas] {
    |[fill=encabezado]| \textbf{Usuarios} \\
    id : BIGINT (PK) \\
    activo : BOOLEAN \\
    fecha\_expiracion\_codigo : TIMESTAMP \\
    nombre : VARCHAR(255) \\
    correo\_institucional : VARCHAR(255) \\
    boleta : VARCHAR(255) \\
    password : VARCHAR(255) \\
    unidad\_academica : VARCHAR(255) \\
    correo\_personal : VARCHAR(255) \\
    telefono\_celular : VARCHAR(255) \\
    nombre\_tutor : VARCHAR(255) \\
    parentesco\_tutor : VARCHAR(255) \\
    telefono\_tutor : VARCHAR(255) \\
    codigo\_recuperacion : VARCHAR(255) \\
};
\matrix (gNotificaciones) [tablaUML, right=4.5cm of gQuejas.north east, anchor=north west] {
    |[fill=encabezado]| \textbf{Notificaciones} \\
    id : BIGINT (PK) \\
    usuario\_id : BIGINT (FK) \\
    mensaje : TEXT \\
    fecha\_creacion : TIMESTAMP \\
    fecha\_expiracion : TIMESTAMP \\
    leida : BOOLEAN \\
    titulo : VARCHAR(255) \\
    tipo : VARCHAR(50) \\
    folio\_queja : VARCHAR(255) \\
};
\matrix (gEvidencias) [tablaUML, left=4cm of gQuejas.south west, anchor=south east] {
    |[fill=encabezado]| \textbf{Evidencias} \\
    id : BIGINT (PK) \\
    queja\_id : BIGINT (FK) \\
    tamano : BIGINT \\
    nombre\_archivo : VARCHAR(255) \\
    url\_almacenamiento : VARCHAR(255) \\
};
\matrix (gTutores) [tablaUML, below=2.8cm of gQuejas] {
    |[fill=encabezado]| \textbf{Tutores} \\
    id : BIGINT (PK) \\
    queja\_id : BIGINT (FK) \\
    nombre : VARCHAR(255) \\
    primer\_apellido : VARCHAR(255) \\
    segundo\_apellido : VARCHAR(255) \\
    parentesco : VARCHAR(255) \\
    correo : VARCHAR(255) \\
    telefono : VARCHAR(255) \\
    url\_identificacion : VARCHAR(255) \\
};
\matrix (gHistorial) [tablaUML, right=4.5cm of gQuejas.south east, anchor=south west] {
    |[fill=encabezado]| \textbf{Historial} \\
    id : BIGINT (PK) \\
    queja\_id : BIGINT (FK) \\
    usuario\_origen : BIGINT (FK) \\
    usuario\_destino : BIGINT (FK) \\
    fecha\_cambio : TIMESTAMP \\
    estado\_anterior : VARCHAR(50) \\
    estado\_nuevo : VARCHAR(50) \\
    observaciones : TEXT \\
};
\draw[lineaRel] (gUsuarios.south) -- node[right, font=\tiny] {1} node[near end, left, font=\tiny] {0..*} (gQuejas.north);
\draw[lineaRel] (gUsuarios.east) -- ++(1,0) |- node[above, font=\tiny] {1} node[near end, above, font=\tiny] {0..*} (gNotificaciones.west);
\draw[lineaRel] (gQuejas.west) -- ++(-1,0) |- node[above, font=\tiny] {1} node[near end, above, font=\tiny] {0..*} (gEvidencias.east);
\draw[lineaRel] (gQuejas.south) -- node[right, font=\tiny] {1} node[near end, left, font=\tiny] {0..1} (gTutores.north);
\draw[lineaRel] (gQuejas.east) -- ++(0.5,0) |- node[above, font=\tiny] {1} node[near end, above, font=\tiny] {0..*} (gHistorial.west);
\end{tikzpicture}
}
\caption{Diagrama Entidad-Relación Global}
\end{figure}

\subsection{Diagramas por módulo}

\subsubsection{Usuarios y Notificaciones}
\begin{figure}[H]
\centering
\begin{tikzpicture}[node distance=3cm, auto]
\matrix (usuarios) [tablaUML] {
    |[fill=encabezado]| \textbf{Usuarios} \\
    id : BIGINT (PK) \\
    nombre : VARCHAR(255) \\
    correo\_institucional : VARCHAR(255) \\
    password : VARCHAR(255) \\
    activo : BOOLEAN \\
};
\matrix (notificaciones) [tablaUML, right=5.5cm of usuarios] {
    |[fill=encabezado]| \textbf{Notificaciones} \\
    id : BIGINT (PK) \\
    usuario\_id : BIGINT (FK) \\
    mensaje : TEXT \\
    fecha\_creacion : TIMESTAMP \\
    leida : BOOLEAN \\
};
\draw[lineaRel] (usuarios.east) -- node[above, font=\tiny] {1} node[near end, above, font=\tiny] {0..*} (notificaciones.west);
\end{tikzpicture}
\caption{DER - Usuarios y Notificaciones}
\end{figure}

\subsubsection{Quejas y Tutores}
\begin{figure}[H]
\centering
\begin{tikzpicture}[node distance=2.5cm, auto]
\matrix (quejas) [tablaUML] {
    |[fill=encabezado]| \textbf{Quejas} \\
    id : BIGINT (PK) \\
    folio : VARCHAR(255) \\
    asunto : VARCHAR(255) \\
    estatus : VARCHAR(50) \\
    fecha\_registro : TIMESTAMP \\
};
\matrix (tutores) [tablaUML, right=5.5cm of quejas] {
    |[fill=encabezado]| \textbf{Tutores} \\
    id : BIGINT (PK) \\
    queja\_id : BIGINT (FK) \\
    nombre : VARCHAR(255) \\
    parentesco : VARCHAR(255) \\
};
\draw[lineaRel] (quejas.east) -- node[above, font=\tiny] {1} node[near end, above, font=\tiny] {0..1} (tutores.west);
\end{tikzpicture}
\caption{DER - Quejas y Tutores}
\end{figure}

\subsubsection{Evidencias}
\begin{figure}[H]
\centering
\begin{tikzpicture}[node distance=3cm, auto]
\matrix (quejas) [tablaUML] {
    |[fill=encabezado]| \textbf{Quejas} \\
    id : BIGINT (PK) \\
    folio : VARCHAR(255) \\
};
\matrix (evidencias) [tablaUML, right=5.5cm of quejas] {
    |[fill=encabezado]| \textbf{Evidencias} \\
    id : BIGINT (PK) \\
    queja\_id : BIGINT (FK) \\
    nombre\_archivo : VARCHAR(255) \\
    url\_almacenamiento : VARCHAR(255) \\
};
\draw[lineaRel] (quejas.east) -- node[above, font=\tiny] {1} node[near end, above, font=\tiny] {0..*} (evidencias.west);
\end{tikzpicture}
\caption{DER - Evidencias}
\end{figure}

\section{Diccionario de Datos}

\subsection{Usuarios}
\begin{longtable}{|p{4.2cm}|p{3.2cm}|p{2.2cm}|p{4.5cm}|}
\hline
\rowcolor{lightgray} \textbf{Campo} & \textbf{Tipo} & \textbf{Nulo} & \textbf{Descripción} \\
\hline
id & BIGINT & NO & Llave primaria autoincremental \\
\hline
activo & BOOLEAN & NO & Estado de la cuenta \\
\hline
fecha\_expiracion\_codigo & TIMESTAMP & SI & Expiración del token de recuperación \\
\hline
nombre & VARCHAR(255) & NO & Nombre(s) y apellidos \\
\hline
correo\_institucional & VARCHAR(255) & NO & Correo único (@ipn.mx) \\
\hline
boleta & VARCHAR(255) & NO & Número de boleta o empleado \\
\hline
password & VARCHAR(255) & NO & Hash BCrypt \\
\hline
unidad\_academica & VARCHAR(255) & SI & Escuela o unidad de adscripción \\
\hline
correo\_personal & VARCHAR(255) & SI & Correo alternativo \\
\hline
telefono\_celular & VARCHAR(255) & SI & Teléfono de contacto \\
\hline
nombre\_tutor & VARCHAR(255) & SI & Nombre del tutor (si aplica) \\
\hline
parentesco\_tutor & VARCHAR(255) & SI & Parentesco con el tutor \\
\hline
telefono\_tutor & VARCHAR(255) & SI & Teléfono del tutor \\
\hline
codigo\_recuperacion & VARCHAR(255) & SI & Código para recuperación \\
\hline
\end{longtable}

\subsection{Quejas}
\begin{longtable}{|p{4.2cm}|p{3.2cm}|p{2.2cm}|p{4.5cm}|}
\hline
\rowcolor{lightgray} \textbf{Campo} & \textbf{Tipo} & \textbf{Nulo} & \textbf{Descripción} \\
\hline
id & BIGINT & NO & Llave primaria autoincremental \\
\hline
usuario\_id & BIGINT & NO & FK Usuarios(id) \\
\hline
descripcion & TEXT & NO & Narrativa de los hechos \\
\hline
fecha\_hechos & TIMESTAMP & NO & Fecha de los hechos \\
\hline
fecha\_registro & TIMESTAMP & NO & Fecha de registro \\
\hline
folio & VARCHAR(255) & NO & Folio único de la queja \\
\hline
asunto & VARCHAR(255) & NO & Título/resumen \\
\hline
unidad\_donde\_ocurrio & VARCHAR(255) & SI & Unidad académica \\
\hline
estatus & VARCHAR(50) & NO & Estado actual \\
\hline
nombre\_denunciado & VARCHAR(255) & SI & Persona denunciada \\
\hline
correo\_quejoso & VARCHAR(255) & NO & Correo del quejoso \\
\hline
nombre\_quejoso & VARCHAR(255) & NO & Nombre del quejoso \\
\hline
identificacion\_institucional & VARCHAR(255) & SI & Número de identificación \\
\hline
tipo\_identificacion & VARCHAR(255) & SI & Tipo (boleta, empleado) \\
\hline
fecha\_nacimiento & VARCHAR(255) & SI & Fecha de nacimiento \\
\hline
\end{longtable}

\subsection{Evidencias}
\begin{longtable}{|p{4.2cm}|p{3.2cm}|p{2.2cm}|p{4.5cm}|}
\hline
\rowcolor{lightgray} \textbf{Campo} & \textbf{Tipo} & \textbf{Nulo} & \textbf{Descripción} \\
\hline
id & BIGINT & NO & Llave primaria \\
\hline
queja\_id & BIGINT & NO & FK Quejas(id) \\
\hline
tamano & BIGINT & NO & Tamaño en bytes \\
\hline
nombre\_archivo & VARCHAR(255) & NO & Nombre original \\
\hline
url\_almacenamiento & VARCHAR(255) & NO & Ruta en storage \\
\hline
\end{longtable}

\subsection{Notificaciones}
\begin{longtable}{|p{4.2cm}|p{3.2cm}|p{2.2cm}|p{4.5cm}|}
\hline
\rowcolor{lightgray} \textbf{Campo} & \textbf{Tipo} & \textbf{Nulo} & \textbf{Descripción} \\
\hline
id & BIGINT & NO & Llave primaria \\
\hline
usuario\_id & BIGINT & NO & FK Usuarios(id) \\
\hline
mensaje & TEXT & NO & Contenido \\
\hline
fecha\_creacion & TIMESTAMP & NO & Fecha de creación \\
\hline
fecha\_expiracion & TIMESTAMP & SI & Fecha de expiración \\
\hline
leida & BOOLEAN & NO & Leída o no \\
\hline
titulo & VARCHAR(255) & NO & Título \\
\hline
tipo & VARCHAR(50) & NO & Tipo (INFO, ALERTA, URGENTE) \\
\hline
folio\_queja & VARCHAR(255) & SI & Folio relacionado \\
\hline
\end{longtable}

\subsection{Historial}
\begin{longtable}{|p{4.2cm}|p{3.2cm}|p{2.2cm}|p{4.5cm}|}
\hline
\rowcolor{lightgray} \textbf{Campo} & \textbf{Tipo} & \textbf{Nulo} & \textbf{Descripción} \\
\hline
id & BIGINT & NO & Llave primaria \\
\hline
queja\_id & BIGINT & NO & FK Quejas(id) \\
\hline
usuario\_origen & BIGINT & SI & FK Usuarios(id) que envía \\
\hline
usuario\_destino & BIGINT & SI & FK Usuarios(id) que recibe \\
\hline
fecha\_cambio & TIMESTAMP & NO & Fecha del cambio \\
\hline
estado\_anterior & VARCHAR(50) & SI & Estado antes \\
\hline
estado\_nuevo & VARCHAR(50) & NO & Estado después \\
\hline
observaciones & TEXT & SI & Comentarios \\
\hline
\end{longtable}

\subsection{Tutores}
\begin{longtable}{|p{4.2cm}|p{3.2cm}|p{2.2cm}|p{4.5cm}|}
\hline
\rowcolor{lightgray} \textbf{Campo} & \textbf{Tipo} & \textbf{Nulo} & \textbf{Descripción} \\
\hline
id & BIGINT & NO & Llave primaria \\
\hline
queja\_id & BIGINT & NO & FK Quejas(id) \\
\hline
nombre & VARCHAR(255) & NO & Nombre del tutor \\
\hline
primer\_apellido & VARCHAR(255) & NO & Primer apellido \\
\hline
segundo\_apellido & VARCHAR(255) & SI & Segundo apellido \\
\hline
parentesco & VARCHAR(255) & NO & Parentesco \\
\hline
correo & VARCHAR(255) & SI & Correo electrónico \\
\hline
telefono & VARCHAR(255) & NO & Teléfono \\
\hline
url\_identificacion & VARCHAR(255) & SI & URL de identificación \\
\hline
\end{longtable}

\section{Relaciones e Integridad}

\subsection{Llaves primarias}
\begin{longtable}{|p{4cm}|p{4cm}|p{6cm}|}
\hline
\rowcolor{lightgray} \textbf{Tabla} & \textbf{Llave primaria} & \textbf{Tipo} \\
\hline
Usuarios & id & BIGINT (SERIAL) \\
\hline
Quejas & id & BIGINT (SERIAL) \\
\hline
Notificaciones & id & BIGINT (SERIAL) \\
\hline
Evidencias & id & BIGINT (SERIAL) \\
\hline
Tutores & id & BIGINT (SERIAL) \\
\hline
Historial & id & BIGINT (SERIAL) \\
\hline
\end{longtable}

\subsection{Llaves foráneas}
\begin{longtable}{|p{4cm}|p{4cm}|p{6cm}|}
\hline
\rowcolor{lightgray} \textbf{Tabla} & \textbf{Llave foránea} & \textbf{Referencia} \\
\hline
Quejas & usuario\_id & Usuarios(id) \\
\hline
Notificaciones & usuario\_id & Usuarios(id) \\
\hline
Evidencias & queja\_id & Quejas(id) \\
\hline
Tutores & queja\_id & Quejas(id) \\
\hline
Historial & queja\_id & Quejas(id) \\
\hline
Historial & usuario\_origen & Usuarios(id) \\
\hline
Historial & usuario\_destino & Usuarios(id) \\
\hline
\end{longtable}

\subsection{Cardinalidades}
\begin{longtable}{|p{4cm}|p{3cm}|p{3cm}|p{4cm}|}
\hline
\rowcolor{lightgray} \textbf{Relación} & \textbf{Origen} & \textbf{Destino} & \textbf{Descripción} \\
\hline
Usuarios -> Quejas & 1 & 0..* & Un usuario puede interponer muchas quejas \\
\hline
Usuarios -> Notificaciones & 1 & 0..* & Un usuario puede recibir muchas notificaciones \\
\hline
Quejas -> Evidencias & 1 & 0..* & Una queja puede tener muchas evidencias \\
\hline
Quejas -> Tutores & 1 & 0..1 & Una queja puede tener 0 o 1 tutor (si es menor) \\
\hline
Quejas -> Historial & 1 & 0..* & Una queja puede tener muchos registros \\
\hline
\end{longtable}

\subsection{Reglas de integridad referencial}
\begin{itemize}
\item CASCADE: Al eliminar un usuario, se eliminan sus notificaciones
\item RESTRICT: No se puede eliminar una queja con evidencias asociadas
\item SET NULL: Al eliminar un usuario, el campo origen/destino en Historial se establece NULL
\item NO ACTION: Las quejas no pueden quedar sin usuario\_id válido
\end{itemize}

\section{Reglas de Negocio}

\subsection{Estados del sistema}
\begin{longtable}{|p{4cm}|p{9cm}|}
\hline
\rowcolor{lightgray} \textbf{Estado} & \textbf{Descripción} \\
\hline
RECIBIDA & La queja ha sido registrada inicialmente \\
\hline
EN\_VALIDACION & Recepción está validando documentos \\
\hline
VALIDADA & Documentación completa y válida \\
\hline
EN\_ANALISIS & Primer Contacto está analizando \\
\hline
CANALIZADA & Titular la ha canalizado a Subdefensoría \\
\hline
EN\_INVESTIGACION & Subdefensoría está investigando \\
\hline
RESUELTA & Expediente concluido \\
\hline
RECHAZADA & Queja no procede \\
\hline
CERRADA & Expediente cerrado administrativamente \\
\hline
\end{longtable}

\subsection{Validaciones}
\begin{enumerate}
\item Correo institucional único: No pueden existir dos usuarios con el mismo correo\_institucional
\item Boleta única: No pueden existir dos usuarios con la misma boleta
\item Folio único: Cada queja tiene un folio único generado automáticamente
\item Mayoría de edad: Si fecha\_nacimiento indica que es menor de edad, se requiere un tutor
\item Código recuperación: El código\_recuperacion expira después de fecha\_expiracion\_codigo
\end{enumerate}

\subsection{Restricciones importantes}
\begin{lstlisting}
ALTER TABLE quejas ADD CONSTRAINT check_menor_tutor CHECK (
    (EXTRACT(YEAR FROM AGE(fecha_nacimiento::DATE)) >= 18) OR 
    (EXTRACT(YEAR FROM AGE(fecha_nacimiento::DATE)) < 18 AND 
     EXISTS (SELECT 1 FROM tutores WHERE tutores.queja_id = quejas.id))
);

ALTER TABLE quejas ADD CONSTRAINT check_estatus_validos CHECK (
    estatus IN ('RECIBIDA','EN_VALIDACION','VALIDADA','EN_ANALISIS',
                'CANALIZADA','EN_INVESTIGACION','RESUELTA','RECHAZADA','CERRADA')
);

ALTER TABLE notificaciones ADD CONSTRAINT check_fecha_expiracion CHECK (
    fecha_expiracion IS NULL OR fecha_expiracion > fecha_creacion
);
\end{lstlisting}

\section{Seguridad y Control de Acceso}

\subsection{Roles de usuario}
\begin{longtable}{|p{4cm}|p{9cm}|}
\hline
\rowcolor{lightgray} \textbf{Rol} & \textbf{Descripción} \\
\hline
ROLE\_QUEJOSO & Usuario que solo puede crear quejas y consultar sus propias quejas \\
\hline
ROLE\_RECEPCION & Personal de recepción - valida y clasifica quejas \\
\hline
ROLE\_PRIMER\_CONTACTO & Analista - realiza análisis preliminar \\
\hline
ROLE\_TITULAR & Director - canaliza expedientes \\
\hline
ROLE\_SUBDEFENSOR & Investigador - gestiona investigación \\
\hline
ROLE\_ADMIN & Administrador del sistema \\
\hline
\end{longtable}

\subsection{Permisos}
\begin{longtable}{|p{4cm}|p{3.5cm}|p{6cm}|}
\hline
\rowcolor{lightgray} \textbf{Rol} & \textbf{Tabla} & \textbf{Operaciones permitidas} \\
\hline
QUEJOSO & Quejas & INSERT (solo propias), SELECT (solo propias) \\
\hline
RECEPCION & Quejas & SELECT, UPDATE (estatus) \\
\hline
PRIMER\_CONTACTO & Quejas & SELECT, UPDATE (análisis) \\
\hline
TITULAR & Quejas & SELECT, UPDATE (canalización) \\
\hline
SUBDEFENSOR & Quejas & SELECT, UPDATE (investigación) \\
\hline
ADMIN & Todas & Todos los privilegios \\
\hline
\end{longtable}

\subsection{Protección de datos}
\begin{itemize}
\item Contraseñas: Almacenadas con BCrypt (factor 10)
\item Evidencias: Hash SHA-256 para verificar integridad
\item Datos sensibles: Los campos identificacion\_institucional y nombre\_denunciado tienen acceso restringido
\item Logs de auditoría: Tabla Historial registra todos los cambios de estado
\item TLS/SSL: Toda comunicación encriptada
\item Row Level Security (RLS): Implementado en PostgreSQL para aislamiento por usuario
\end{itemize}

\section{Índices y Optimización}

\subsection{Índices principales}
\begin{lstlisting}
CREATE INDEX idx_usuarios_correo_institucional ON usuarios(correo_institucional);
CREATE INDEX idx_usuarios_boleta ON usuarios(boleta);
CREATE INDEX idx_usuarios_codigo_recuperacion ON usuarios(codigo_recuperacion);
CREATE INDEX idx_quejas_usuario_id ON quejas(usuario_id);
CREATE INDEX idx_quejas_folio ON quejas(folio);
CREATE INDEX idx_quejas_estatus ON quejas(estatus);
CREATE INDEX idx_quejas_fecha_registro ON quejas(fecha_registro);
CREATE INDEX idx_notificaciones_usuario_id ON notificaciones(usuario_id);
CREATE INDEX idx_notificaciones_leida ON notificaciones(leida);
CREATE INDEX idx_evidencias_queja_id ON evidencias(queja_id);
CREATE INDEX idx_tutores_queja_id ON tutores(queja_id);
CREATE INDEX idx_historial_queja_id ON historial(queja_id);
CREATE INDEX idx_historial_fecha_cambio ON historial(fecha_cambio);
\end{lstlisting}

\subsection{Optimización de consultas}
\begin{lstlisting}
SELECT q.id, q.folio, q.asunto, q.estatus, u.nombre as quejoso
FROM quejas q
JOIN usuarios u ON q.usuario_id = u.id
WHERE q.estatus IN ('RECIBIDA', 'EN_VALIDACION')
ORDER BY q.fecha_registro DESC;

SELECT * FROM quejas WHERE folio = $1;

SELECT * FROM notificaciones 
WHERE usuario_id = $1 AND leida = false 
ORDER BY fecha_creacion DESC;
\end{lstlisting}

\section{Scripts SQL}

\subsection{Creación de tablas}
\begin{lstlisting}
CREATE TABLE usuarios (
    id BIGSERIAL PRIMARY KEY,
    activo BOOLEAN DEFAULT TRUE,
    fecha_expiracion_codigo TIMESTAMP,
    nombre VARCHAR(255) NOT NULL,
    correo_institucional VARCHAR(255) UNIQUE NOT NULL,
    boleta VARCHAR(255) UNIQUE NOT NULL,
    password VARCHAR(255) NOT NULL,
    unidad_academica VARCHAR(255),
    correo_personal VARCHAR(255),
    telefono_celular VARCHAR(255),
    nombre_tutor VARCHAR(255),
    parentesco_tutor VARCHAR(255),
    telefono_tutor VARCHAR(255),
    codigo_recuperacion VARCHAR(255)
);

CREATE TABLE quejas (
    id BIGSERIAL PRIMARY KEY,
    usuario_id BIGINT NOT NULL REFERENCES usuarios(id),
    descripcion TEXT NOT NULL,
    fecha_hechos TIMESTAMP NOT NULL,
    fecha_registro TIMESTAMP DEFAULT CURRENT_TIMESTAMP,
    folio VARCHAR(255) UNIQUE NOT NULL,
    asunto VARCHAR(255) NOT NULL,
    unidad_donde_ocurrio VARCHAR(255),
    estatus VARCHAR(50) NOT NULL,
    nombre_denunciado VARCHAR(255),
    correo_quejoso VARCHAR(255) NOT NULL,
    nombre_quejoso VARCHAR(255) NOT NULL,
    identificacion_institucional VARCHAR(255),
    tipo_identificacion VARCHAR(255),
    fecha_nacimiento VARCHAR(255)
);

CREATE TABLE evidencias (
    id BIGSERIAL PRIMARY KEY,
    queja_id BIGINT NOT NULL REFERENCES quejas(id) ON DELETE CASCADE,
    tamano BIGINT NOT NULL,
    nombre_archivo VARCHAR(255) NOT NULL,
    url_almacenamiento VARCHAR(255) NOT NULL
);

CREATE TABLE notificaciones (
    id BIGSERIAL PRIMARY KEY,
    usuario_id BIGINT NOT NULL REFERENCES usuarios(id) ON DELETE CASCADE,
    mensaje TEXT NOT NULL,
    fecha_creacion TIMESTAMP DEFAULT CURRENT_TIMESTAMP,
    fecha_expiracion TIMESTAMP,
    leida BOOLEAN DEFAULT FALSE,
    titulo VARCHAR(255) NOT NULL,
    tipo VARCHAR(50) NOT NULL,
    folio_queja VARCHAR(255)
);

CREATE TABLE tutores (
    id BIGSERIAL PRIMARY KEY,
    queja_id BIGINT NOT NULL REFERENCES quejas(id) ON DELETE CASCADE,
    nombre VARCHAR(255) NOT NULL,
    primer_apellido VARCHAR(255) NOT NULL,
    segundo_apellido VARCHAR(255),
    parentesco VARCHAR(255) NOT NULL,
    correo VARCHAR(255),
    telefono VARCHAR(255) NOT NULL,
    url_identificacion VARCHAR(255)
);

CREATE TABLE historial (
    id BIGSERIAL PRIMARY KEY,
    queja_id BIGINT NOT NULL REFERENCES quejas(id) ON DELETE CASCADE,
    usuario_origen BIGINT REFERENCES usuarios(id) ON DELETE SET NULL,
    usuario_destino BIGINT REFERENCES usuarios(id) ON DELETE SET NULL,
    fecha_cambio TIMESTAMP DEFAULT CURRENT_TIMESTAMP,
    estado_anterior VARCHAR(50),
    estado_nuevo VARCHAR(50) NOT NULL,
    observaciones TEXT
);
\end{lstlisting}

\subsection{Índices (ejecutar después de creación)}
\begin{lstlisting}
CREATE INDEX idx_usuarios_correo_institucional ON usuarios(correo_institucional);
CREATE INDEX idx_usuarios_boleta ON usuarios(boleta);
CREATE INDEX idx_quejas_usuario_id ON quejas(usuario_id);
CREATE INDEX idx_quejas_folio ON quejas(folio);
CREATE INDEX idx_quejas_estatus ON quejas(estatus);
CREATE INDEX idx_notificaciones_usuario_id ON notificaciones(usuario_id);
CREATE INDEX idx_evidencias_queja_id ON evidencias(queja_id);
CREATE INDEX idx_tutores_queja_id ON tutores(queja_id);
CREATE INDEX idx_historial_queja_id ON historial(queja_id);
\end{lstlisting}

\subsection{Constraints adicionales}
\begin{lstlisting}
ALTER TABLE quejas ADD CONSTRAINT check_estatus_validos CHECK (
    estatus IN ('RECIBIDA','EN_VALIDACION','VALIDADA','EN_ANALISIS',
                'CANALIZADA','EN_INVESTIGACION','RESUELTA','RECHAZADA','CERRADA')
);
ALTER TABLE notificaciones ADD CONSTRAINT check_tipo_validos CHECK (
    tipo IN ('INFO','ALERTA','URGENTE')
);
ALTER TABLE notificaciones ADD CONSTRAINT check_fecha_expiracion CHECK (
    fecha_expiracion IS NULL OR fecha_expiracion > fecha_creacion
);
\end{lstlisting}

\subsection{Datos de prueba (opcional)}
\begin{lstlisting}
INSERT INTO usuarios (nombre, correo_institucional, boleta, password, activo)
VALUES ('Juan Perez Garcia', 'jperez@ipn.mx', '2024A12345', '$2a$10$...', true);

INSERT INTO quejas (usuario_id, descripcion, fecha_hechos, folio, asunto, 
                    estatus, correo_quejoso, nombre_quejoso)
VALUES (1, 'Descripcion de prueba', '2024-01-15 10:00:00', 
        'DDP-2024-0001', 'Problema con servicio', 'RECIBIDA', 
        'jperez@ipn.mx', 'Juan Perez Garcia');
\end{lstlisting}

\end{document}